\lecture{Why FEM, and the two-point BVP}{Lecture 1 --- Monday, August 24, 2026}{Math 6660}

\begin{roadmap}[Lecture 1]
\roadq{Why does FEM not just solve the differential equation?}
\begin{roadlist}
\item The model problem for the whole course is $-(au')'=f$. Read it as an
  \textbf{axial bar}: $a=EA$ is the stiffness, $f$ the distributed load, $u$ the
  displacement. Everything later is this equation or a relative of it.
\item FEM never attacks the \emph{strong} form (the ODE holding at every point).
  It rewrites the problem as one \emph{integral} identity, tested against
  arbitrary virtual displacements. That is the weak form.
\item The operator turns out to be \textbf{self-adjoint} --- which is
  Betti--Maxwell reciprocity wearing a different hat.
\item The space $\Vsp$ of admissible displacements is defined by
  \textbf{finite strain energy}, not by smoothness. That choice is what lets
  kinked, piecewise-linear elements be legal functions.
\end{roadlist}
\payoff{Knowing why every element derivation you have ever seen starts from
virtual work rather than from the governing PDE.}
\end{roadmap}

% ==================================================================
\subsection*{Context}

Where this course sits relative to the rest of the numerical sequence:

\begin{center}
\begin{tabular}{@{}ll@{}}
\toprule
\textbf{Prerequisites} & \textbf{The discretisation courses}\\
\midrule
5610: interpolation, integration, ODEs & \textbf{6660}: FEMs \quad(this course)\\
5620: numerical linear algebra          & 6670: FVMs\\
                                        & 5170: FDMs\\
\bottomrule
\end{tabular}
\end{center}

% ==================================================================
\subsection*{Why FEM?}

\paragraph{(1) Theoretical.}
The original problem is continuous.
\begin{itemize}[leftmargin=1.6em,itemsep=1pt]
  \item \textbf{FDM}: approximate the solution at a discrete number of points.
  \item \textbf{FEM}: keep the original continuous problem (the PDE), but
        replace (approximate) the \emph{function spaces}.
\end{itemize}
All the theory of PDEs carries over to FEM.

\begin{intuition}[FDM discretises the equation; FEM discretises the space]
This is the single most important sentence of the lecture, and it is the
reason the rest of the course looks like functional analysis rather than
calculus.

FDM replaces $\dd/\dx$ by a difference quotient --- it attacks the
\emph{operator}, and what survives is a list of nodal values with no meaning in
between. FEM leaves the equation alone and instead shrinks the \emph{set of
candidate shapes} from ``all displacement fields with finite energy'' to ``all
piecewise polynomials on this mesh.'' The answer is still a function defined
everywhere; you can differentiate it, integrate it, and ask for its value
between nodes.

That is why an FE result gives you a stress field over the whole part while a
finite-difference result gives you numbers at grid points, and it is why FEM
inherits the PDE's theory: the problem was never changed, only the space it is
posed over.
\end{intuition}

\paragraph{(2) Practical reasons.}
\begin{itemize}[leftmargin=1.6em,itemsep=1pt]
  \item Can be formulated on an unstructured grid.
  \item Lots of flexibility in the choice of function spaces.
\end{itemize}

\begin{mechanics}
Both bullets are why FEM won in engineering practice and FDM did not. An
unstructured grid is what lets you mesh a bracket with a fillet and three bolt
holes; a structured finite-difference grid cannot follow that boundary.
``Flexibility in the choice of function spaces'' is element selection --- linear
vs.\ quadratic, Lagrange vs.\ Hermite, beam vs.\ shell vs.\ solid. Choosing an
element type \emph{is} choosing $\Vh$.
\end{mechanics}

\paragraph{(3) Note.}
Much of the theory of FEMs is based on \emph{elliptic} PDE theory. Parabolic
PDEs are OK. Hyperbolic PDEs are more challenging.

\begin{intuition}[why elliptic problems are the friendly case]
Elliptic $=$ steady-state and diffusive: static stress analysis, steady heat
conduction, incompressible potential flow. Information propagates
instantaneously in all directions, the solution is as smooth as the data
allows, and --- the part that matters for this course --- the problem has an
associated \emph{energy} that the solution minimises. That energy is what
becomes the bilinear form $a(u,v)$, and its positivity is what makes every
existence proof work.

Hyperbolic $=$ wave-like: information travels along characteristics at finite
speed, discontinuities (shocks, wavefronts) propagate rather than smooth out,
and there is no energy functional to minimise. Standard Galerkin FEM has no
mechanism to respect a direction of propagation, which is why it produces
spurious oscillations on advection-dominated problems and why that case needs
separate machinery (SUPG, discontinuous Galerkin). If you have run an explicit
crash or blast simulation and wondered why it uses such different technology
from a static solve, this is why.
\end{intuition}

% ==================================================================
\subsection*{The two-point BVP}

The 1D version of Poisson's equation. On the fixed domain $0<x<1$:
\begin{equation}\label{eq:bvp}
  -\bigl(a(x)\,u'\bigr)' = f(x),
  \qquad u(0)=u(1)=0 .
\end{equation}

\begin{center}
\begin{tikzpicture}[scale=1.0]
  % ---- the bar, loaded AXIALLY ----
  \draw[thick,fill=IntBack] (0,-0.22) rectangle (6,0.22);
  % supports (hatched walls)
  \draw[thick] (0,-0.6) -- (0,0.6);
  \foreach \y in {-0.55,-0.35,...,0.45} \draw (0,\y) -- (-0.22,\y+0.13);
  \draw[thick] (6,-0.6) -- (6,0.6);
  \foreach \y in {-0.55,-0.35,...,0.45} \draw (6,\y) -- (6.22,\y+0.13);
  % axial distributed load: arrows ALONG the bar
  \foreach \x in {0.7,1.5,...,5.1}
    \draw[-{Latex[length=2mm]},MEFrame,thick] (\x,0) -- (\x+0.55,0);
  \node[MEFrame] at (3.0,0.62) {$f(x)$ \ \small(axial)};
  \node at (0,-0.85) {\scriptsize$0$};
  \node at (6,-0.85) {\scriptsize$1$};
  % ---- u(x) plotted below, NOT as a sag ----
  \begin{scope}[shift={(0,-2.6)}]
    \draw[->,gray!70] (-0.2,0) -- (6.6,0) node[right,font=\small,black] {$x$};
    \draw[->,gray!70] (0,-0.15) -- (0,1.25);
    \node[right,font=\small] at (0.12,1.05) {$u(x)$};
    \draw[very thick,IntFrame] (0,0) .. controls (2,0.95) and (4,0.95) .. (6,0);
    \foreach \x/\l in {0/0, 6/1}{\draw (\x,-0.1)--(\x,0.1); \node[below,font=\scriptsize] at (\x,-0.12) {$\l$};}
  \end{scope}
\end{tikzpicture}
\end{center}

\begin{pitfall}[do not draw this as a sagging beam]
The obvious picture --- downward arrows and a drooping curve --- is
\textbf{wrong}, and the lecturer flags it directly: \emph{``misleading figure;
try to think of it as a heat transfer problem rather than a structural
mechanics problem.''}

Here $u$ is \emph{axial} displacement and $f$ is an \emph{axial} load, so the
load arrows point along the bar, and the curve below is a \emph{plot of $u$
against $x$}, not the shape the bar takes. A sagging picture silently imports
bending, which is a fourth-order problem with a different energy
($\int EI(w'')^{2}$) and different elements --- see the $H^{1}$ vs $H^{2}$
discussion later in this lecture.

The safest reading of \eqref{eq:bvp} is steady heat conduction: $u$
temperature, $a$ conductivity, $f$ a heat source. Nothing there tempts you into
bending. Use the bar for the \emph{algebra} ($au'=N$, $\int a(u')^{2}=$ strain
energy); use heat conduction for the \emph{picture}.
\end{pitfall}

\noindent
\begin{tabular}{@{}r@{\;}c@{\;}l@{\qquad}l@{}}
$u$ & $\defeq$ & ``displacement''          & \\
$f$ & $\defeq$ & ``external load''         & given\\
$a$ & $=$      & ``modulus of elasticity'' & given\\
\end{tabular}

\medskip
\noindent\textbf{Assume:}
\begin{equation}\label{eq:abounds}
  0 < a_{1} \le a(x) \le a_{2} < \infty
  \qquad \forall\, x\in[0,1].
\end{equation}

\begin{mechanics}[the model problem is an axially loaded bar]
The lecturer's own naming makes this explicit, so take it literally. Let
$a(x)=E(x)A(x)$. Then $a u'$ is the internal axial force $N(x) = EA\,\eps$, and
\eqref{eq:bvp} reads
\[
  -\frac{\dd N}{\dx} = f
  \qquad\text{i.e.}\qquad
  \frac{\dd N}{\dx} + f = 0,
\]
which is nothing but axial equilibrium of a differential slice of bar. The
boundary condition $u(0)=u(1)=0$ pins both ends.

So \eqref{eq:bvp} is: \emph{a bar, fixed at both ends, under distributed axial
load $f$, made of a material whose stiffness $a$ may vary along its length.}
Every abstract statement in the next four lectures can be checked against this
picture.

Read \eqref{eq:abounds} the same way: $a_1>0$ says the bar is nowhere made of
nothing (no zero-stiffness segment, or it would separate freely), and
$a_2<\infty$ says it is nowhere rigid. These two constants will reappear later
as the coercivity and continuity constants, and their ratio $a_2/a_1$ is the
material contrast that controls how bad the error bound gets.
\end{mechanics}

\todoq{Page 3 of my notes writes the boundary condition as $u(0)=u(1)=1$ at the
top, but every subsequent line --- and all of Lecture 2 --- uses
$u(0)=u(1)=0$. I have transcribed it as $0$. Confirm it was a slip.}

% ==================================================================
\subsection*{The operator is self-adjoint}

Note that
\begin{equation}\label{eq:Ldef}
  \mathcal{L}(u) = -(a u')'
  \qquad\text{with}\qquad u(0)=u(1)=0
\end{equation}
is a linear differential operator, and $\mathcal{L}$ is \emph{self-adjoint}.

To see it, integrate by parts twice:
\begin{align}
  \int_{0}^{1} v\,\mathcal{L}(u)\dx
    &= -\int_{0}^{1} v\,(au')'\dx \notag\\[2pt]
    &= \int_{0}^{1} v'\,a\,u'\dx - \Bigl[v\,a\,u'\Bigr]_{x=0}^{x=1}.
    \label{eq:sa1}
\end{align}
Assume $u(0)=u(1)=0$ and $v(0)=v(1)=0$; the boundary term drops:
\begin{align}
  &= \int_{0}^{1} v'\,a\,u'\dx \notag\\[2pt]
  &= -\int_{0}^{1} u\,(a v')'\dx + \Bigl[u\,a\,v'\Bigr]_{x=0}^{x=1} \notag\\[2pt]
  &= -\int_{0}^{1} u\,(a v')'\dx
   = \int_{0}^{1} u\,\mathcal{L}(v)\dx .
\end{align}

\begin{equation}\label{eq:selfadjoint}
  \boxed{\;\int_{0}^{1} v\,\mathcal{L}(u)\dx = \int_{0}^{1} u\,\mathcal{L}(v)\dx\;}
  \qquad v = \text{test function}
\end{equation}
This is what it means to be self-adjoint.

\todoq{My notes write line \eqref{eq:sa1} as
$\int v\,\mathcal{L}(u) = \int v\,(au')'$, without the minus sign, even though
$\mathcal{L}(u)=-(au')'$ two lines above. The very next line \emph{does} carry
the sign correctly, so it is a copying slip and I have restored the minus here.}

\begin{intuition}[self-adjointness is Betti's reciprocal theorem]
Equation~\eqref{eq:selfadjoint} is not a piece of abstract bookkeeping --- it is
a theorem you were taught in structures, under a different name.

Read $\mathcal{L}(u)$ as ``the load distribution that produces displacement
field $u$.'' Then the left side is \emph{the work done by load system $u$ moving
through the displacements of system $v$}, and the right side is the work done by
load system $v$ moving through the displacements of $u$. Equation
\eqref{eq:selfadjoint} says these are equal.

That is \textbf{Betti--Maxwell reciprocity}: push at $A$ and measure the
deflection at $B$; you get the same number as pushing at $B$ and measuring at
$A$. It is also exactly why the flexibility and stiffness matrices are
symmetric.
\end{intuition}

\medskip
\noindent
\textbf{Note.} A self-adjoint differential operator $\mathcal{L}$ is analogous
to a symmetric matrix.

\begin{dictionary}
$\mathcal{L}$ self-adjoint & $\Kmat = \Kmat^{T}$ \\
$\int v\,\mathcal{L}(u) = \int u\,\mathcal{L}(v)$ & $\bm{v}^{T}\Kmat\bm{u} = \bm{u}^{T}\Kmat\bm{v}$ \\
the boundary term $[v\,a\,u']_{0}^{1}$ & end reactions doing work; killed by the pinned BCs \\
$a\,u'$ & internal axial force $N = EA\,\eps$ \\
\end{dictionary}

\begin{pitfall}[the boundary term is where BCs actually enter]
Watch what happened in \eqref{eq:sa1}: the essential boundary condition was not
imposed on the equation, it was used to \emph{delete the boundary term}. This
is the mechanism for the whole subject.

Kill the term by forcing $v=0$ at the end $\Rightarrow$ that end carries an
\emph{essential} (Dirichlet, prescribed-displacement) condition, and it must be
built into the space. Kill it instead by prescribing $a u' = $ given
$\Rightarrow$ that end carries a \emph{natural} (Neumann, prescribed-traction)
condition, and it appears in the load vector for free, with nothing imposed on
the space.

That is the real content of ``essential vs.\ natural,'' and it is why in a
commercial FE code you constrain displacement DOFs explicitly but simply
\emph{apply} a traction and let the formulation absorb it.
\end{pitfall}

% ==================================================================
\subsection*{Invertibility and regularity}

\noindent\textbf{Note.} Consider
\[
  \mathcal{L}(w)=0,\qquad w(0)=w(1)=0 .
\]
The only solution that satisfies $\mathcal{L}$ and the BCs is $w=0$. Therefore
$\mathcal{L}$ is ``invertible'': for given $a$ and $f$ there exists a unique
$u$.

\begin{intuition}[trivial null space $=$ no rigid-body modes]
$\mathcal{L}(w)=0$ with $w=0$ at both ends asks: \emph{is there a nonzero
displacement field that requires no load?} For a bar pinned at both ends,
obviously not --- any nonzero shape stretches some fibre, which needs force.
Release one end and the answer changes: $w = \text{const}$ becomes a free
translation, the null space is nontrivial, and the problem has no unique
solution.

In matrix terms: $\Kmat\bm{w}=\bm{0}\Rightarrow\bm{w}=\bm{0}$ means
$\Kmat$ is nonsingular. It is the same check as ``did I constrain enough DOFs
to remove all rigid body motion,'' which is the first thing a solver complains
about when it reports a singular stiffness matrix.
\end{intuition}

\todoq{My notes read ``there exist a unique $u$ and $f$'' --- but $f$ is given
data, not something solved for. I have written ``for given $a$ and $f$ there
exists a unique $u$.'' Check the phrasing.}

\medskip
\noindent\textbf{Note (regularity).} If $0<a<\infty$ and $f$ are sufficiently
smooth --- specifically, if
\[
  a \in \Ck[(0,1)]{1}\quad\text{and}\quad f\in\Ck[(0,1)]{0},
\]
then $u\in\Ck[(0,1)]{2}$.

\begin{intuition}[the solution is two derivatives smoother than the load]
$\mathcal{L}$ takes two derivatives, so solving with $\mathcal{L}^{-1}$ gives
two back. This is the \emph{elliptic regularity} pattern, and it is why the
elliptic case is the friendly one.

It also tells you where it breaks: the hypothesis needs $a\in C^{1}$, and a bar
made of two glued materials has $a$ \emph{jumping} at the interface. Then
$u\notin C^{2}$, the strong form \eqref{eq:bvp} is not satisfied in the
classical sense, and there is no classical solution at all --- even though the
bar plainly has a displacement field. Resolving that is precisely what the
variational formulation is for, and Lecture 2 closes the loop on it.
\end{intuition}

% ==================================================================
\subsection*{Variational formulation}

Multiply our BVP by a \underline{test function} and integrate over the domain
$[0,1]$:
\begin{equation}\label{eq:mult}
  -\int_{0}^{1}(au')'\,v\dx = \int_{0}^{1} \underbrace{f}_{\text{force}} v\dx .
\end{equation}
Then integrate by parts:
\begin{equation}\label{eq:ibp1d}
  \int_{0}^{1} a\,u'v'\dx - \Bigl[a\,u'v\Bigr]_{0}^{1} = \int_{0}^{1} f v\dx .
\end{equation}
Assume $v$ satisfies the BCs:
\[
  v(0)=v(1)=0 .
\]
In particular, define the solution space
\begin{equation}\label{eq:Vdef}
  V = \left\{\, v \;:\;
    \int_{0}^{1} v^{2}\dx < \infty,\quad
    \int_{0}^{1} \bigl[v'\bigr]^{2}\dx < \infty,\quad
    v(0)=v(1)=0 \,\right\}.
\end{equation}

\begin{intuition}[$V$ is ``every shape the bar could actually take'']
Read the three conditions in \eqref{eq:Vdef} physically, in order:
\begin{itemize}[leftmargin=1.4em,itemsep=2pt]
  \item $\int v^{2}<\infty$ --- the displacement is finite in the mean. The bar
        does not fly off to infinity.
  \item $\int (v')^{2}<\infty$ --- \textbf{finite strain energy.} Since
        $v'=\eps$, this integral is (up to the factor $a/2$) the stored elastic
        energy. This is the condition that does all the work.
  \item $v(0)=v(1)=0$ --- kinematically admissible: the shape respects the
        supports.
\end{itemize}
So $V$ is the set of \emph{kinematically admissible displacement fields with
finite strain energy}. Nothing more exotic than that. This space is called
$H^{1}_{0}(0,1)$ once the course introduces Sobolev notation; when you meet that
symbol, substitute this sentence.

\medskip
\noindent\textbf{What each rung of the derivative ladder is.}
Every condition defining $V$ is a statement about one of these, and nothing in
the list goes past the first derivative:
\begin{center}\small
\renewcommand{\arraystretch}{1.15}
\begin{tabular}{@{}llll@{}}
\toprule
\textbf{Object} & \textbf{Mechanics} & \textbf{Units} & \textbf{Where it appears}\\
\midrule
$v$          & displacement            & m       & $\int v^{2}<\infty$; the BCs $v(0)=v(1)=0$\\
$v'$         & strain $\eps$           & ---     & $\int (v')^{2}<\infty$: the energy condition\\
$a v'$       & stress $\sigma$ (times area) & Pa & the flux; continuous across a material joint\\
$(a v')'$    & net force per length    & N/m     & \emph{only} in the strong form $-(av')'=f$\\
\bottomrule
\end{tabular}
\end{center}

\noindent
Because $0<a_{1}\le a\le a_{2}<\infty$, the middle two rows say the same thing:
$a_{1}^{2}\int(v')^{2}\le\int\sigma^{2}\le a_{2}^{2}\int(v')^{2}$, so ``finite
strain energy'', ``$v'$ square integrable'' and ``$\sigma$ square integrable''
are interchangeable. That two-sided bound on $a$ is exactly what keeps them
from coming apart.

\medskip
\noindent\textbf{Why $v''$ is not required.}
Two independent reasons, and it is worth holding both:
\begin{enumerate}[label=(\roman*),leftmargin=1.8em,itemsep=2pt]
  \item \textbf{The energy never contains it.} Strain energy is
        $\tfrac12\int a(v')^{2}$ --- first derivative only. A condition on
        $v''$ would be restricting something the physics does not charge for.
  \item \textbf{Integration by parts removed it.} The strong form demands
        $u''$; \eqref{eq:ibp1d} moved one derivative onto $v$, leaving $u'$ and
        $v'$. After that step $v''$ has left the problem entirely, and $u'$,
        $v'$ appear \emph{only inside integrals} --- never evaluated at a point.
\end{enumerate}

\noindent
Hence what $V$ does \emph{not} require: $v''$ need not exist, and $v'$ need not
exist in the classical sense everywhere. A hat function --- kinked, not
differentiable at its peak --- is in $V$, because that single bad point has zero
length and contributes nothing to $\int(v')^{2}$. That is deliberate: the space
was designed to be exactly big enough to hold the piecewise-linear shapes
finite elements are built from.

The whole distinction comes down to \textbf{kink versus jump}
(Figure~\ref{fig:kinkjump}), and it is worth learning once because the same two
words recur throughout the course.
\end{intuition}

% ------------------------------------------------------------------
\iftranscript\else % --- kink/jump figure: added analysis, not lecture content
\begin{figure}[htbp]
\centering
\begin{tikzpicture}[scale=0.92, every node/.style={font=\small}]

% ---------- column headers ----------
\node[font=\small\bfseries] at (1.6,3.15) {the shape $v(x)$};
\node[font=\small\bfseries] at (6.6,3.15) {its derivative $v'(x)$};

% ================= ROW 1 : KINK =================
\begin{scope}[shift={(0,1.55)}]
  \draw[->,gray] (-0.1,0) -- (3.4,0);
  \draw[very thick,IntFrame] (0.2,0) -- (1.6,1.05) -- (3.0,0);
  \fill[PitFrame] (1.6,1.05) circle (2pt);
  \node[PitFrame,above right,font=\scriptsize] at (1.55,1.0) {corner};
\end{scope}
\begin{scope}[shift={(5,1.55)}]
  \draw[->,gray] (-0.1,0) -- (3.4,0);
  \draw[very thick,IntFrame] (0.2,0.62) -- (1.6,0.62);
  \draw[very thick,IntFrame] (1.6,-0.62) -- (3.0,-0.62);
  \draw[dotted,PitFrame,thick] (1.6,0.62) -- (1.6,-0.62);
  \draw[IntFrame,fill=white] (1.6,0.62) circle (1.7pt);
  \draw[IntFrame,fill=white] (1.6,-0.62) circle (1.7pt);
  \node[right,font=\scriptsize] at (3.05,0.62) {$+1$};
  \node[right,font=\scriptsize] at (3.05,-0.62) {$-1$};
  \node[PitFrame,right,font=\scriptsize] at (1.72,0.05) {finite jump};
\end{scope}
\node[left,font=\small\bfseries,MEFrame] at (-0.35,1.55) {KINK};
\node[right,align=left,font=\scriptsize] at (8.9,1.55)
  {$\displaystyle\int (v')^{2}$ \textbf{finite}\\[1pt] $v\in V$ \ \cmark};

% ================= ROW 2 : JUMP =================
\begin{scope}[shift={(0,-0.75)}]
  \draw[->,gray] (-0.1,0) -- (3.4,0);
  \draw[very thick,IntFrame] (0.2,0) -- (1.1,0);
  \draw[very thick,IntFrame] (1.1,0.95) -- (2.1,0.95);
  \draw[very thick,IntFrame] (2.1,0) -- (3.0,0);
  \draw[dotted,PitFrame,thick] (1.1,0) -- (1.1,0.95);
  \draw[dotted,PitFrame,thick] (2.1,0) -- (2.1,0.95);
  \draw[IntFrame,fill=white] (1.1,0)    circle (1.7pt);
  \draw[IntFrame,fill=white] (2.1,0.95) circle (1.7pt);
  \node[PitFrame,above,font=\scriptsize] at (1.6,1.0) {gap};
\end{scope}
\begin{scope}[shift={(5,-0.75)}]
  \draw[->,gray] (-0.1,0) -- (3.4,0);
  \draw[very thick,-{Latex[length=2.5mm]},PitFrame] (1.1,0) -- (1.1,1.0);
  \draw[very thick,-{Latex[length=2.5mm]},PitFrame] (2.1,0) -- (2.1,-1.0);
  \node[PitFrame,right,font=\scriptsize] at (1.18,0.85) {$+\delta$};
  \node[PitFrame,right,font=\scriptsize] at (2.18,-0.85) {$-\delta$};
  \node[right,font=\scriptsize] at (2.65,0.5) {spikes};
\end{scope}
\node[left,font=\small\bfseries,PitFrame] at (-0.35,-0.75) {JUMP};
\node[right,align=left,font=\scriptsize] at (8.9,-0.75)
  {$\displaystyle\int (v')^{2}=\infty$\\[1pt] $v\notin V$ \ \xmark};

\end{tikzpicture}

\caption{\textbf{Kink versus jump.} A \emph{kink} is a break in the slope: the
shape stays connected, so the strain merely jumps between two finite values and
the energy $\int(v')^{2}$ stays finite. A \emph{jump} is a break in the shape
itself: the strain becomes a delta spike and the energy is infinite.
Mechanically, a kink is a \textbf{material interface} (steel glued to
aluminium: displacement continuous, strain jumps because $E$ jumps) --- ordinary
and everywhere. A jump is a \textbf{broken bar}. $V$ admits the first and
excludes the second, automatically, with no continuity rule written
anywhere.}
\label{fig:kinkjump}
\end{figure}
\fi

\begin{mechanics}[the same rule for beams: why they need rotational DOFs]
The rule generalises: \emph{the highest derivative appearing in the energy
fixes the space, which fixes the continuity the elements must have.} Working it
out for a beam explains an element-selection rule you already follow.

For a beam the derivative ladder is shifted up by one rung, because bending
energy lives in \emph{curvature} rather than in stretch:

\begin{center}\small
\renewcommand{\arraystretch}{1.2}
\begin{tabular}{@{}l l l@{}}
\toprule
 & \textbf{Bar (this course, for now)} & \textbf{Beam (Euler--Bernoulli)}\\
\midrule
field                & $v$ = axial displacement      & $w$ = transverse deflection\\
first derivative     & $v'$ = strain $\eps$          & $w'$ = slope $\theta$\\
second derivative    & ---                            & $w''$ = curvature $\kappa$\\
internal measure     & $\sigma = E v'$                & $M = EI\,w''$ (bending moment)\\
\textbf{energy}      & $\tfrac12\int EA\,(v')^{2}$    & $\tfrac12\int EI\,(w'')^{2}$\\
equilibrium          & $-(EA\,v')' = f$ \ (2nd order) & $(EI\,w'')'' = q$ \ (4th order)\\
space                & $H^{1}$: $\int (v')^{2}<\infty$ & $H^{2}$: $\int (w'')^{2}<\infty$\\
must stay continuous & $v$                            & $w$ \emph{and} $w'$\\
forbidden            & jump in $v$ (a break)          & jump in $w'$ (a hinge)\\
element              & $C^{0}$, DOF $\{w\}$           & $C^{1}$, DOF $\{w,\theta\}$\\
basis                & linear hat (Lagrange)          & cubic Hermite\\
\bottomrule
\end{tabular}
\end{center}

Now run the Figure~\ref{fig:kinkjump} argument one rung higher. For the bar, a
\emph{jump in $v$} gives a delta in $v'$ and infinite energy, so it is
forbidden; a \emph{kink in $v$} is fine. For the beam the energy contains
$w''$, so the thing that must not jump is $w'$:
\begin{center}
\begin{tabular}{@{}l@{\ }c@{\ }l@{\ }c@{\ }l@{}}
bar:  & jump in $v$  & $\Rightarrow$ $\delta$ in $v'$   & $\Rightarrow$ & $\int(v')^{2}=\infty$ \ \xmark\\
beam: & jump in $w'$ & $\Rightarrow$ $\delta$ in $w''$  & $\Rightarrow$ & $\int(w'')^{2}=\infty$ \ \xmark\\
\end{tabular}
\end{center}

\textbf{What a slope jump physically is: a hinge.} A kink in a beam's deflected
shape means the two sides meet at an angle. The curvature there is infinite ---
you have turned through a finite angle over zero length --- so the bending
energy $\tfrac12\int EI\kappa^{2}$ is infinite. A genuine continuous beam
cannot do this. (A real hinge \emph{can}, but only because it has $EI=0$ there,
so it stores no energy; that is a different structure, and you model it by
deliberately releasing the rotational DOF.)

\textbf{The consequence for elements.} Two adjacent bar elements sharing a node
share one number, the displacement --- and that alone makes $v$ continuous, so
$C^{0}$ suffices. Two adjacent beam elements must agree on \emph{both} the
deflection and the rotation at the shared node, or a hinge appears at every
node of your mesh. That is why a beam element carries $\{w_{1},\theta_{1},
w_{2},\theta_{2}\}$ --- four DOFs, hence a cubic ($C^{1}$ Hermite) shape
function --- while a bar element carries only $\{w_{1},w_{2}\}$.

So when a solver gives a beam or shell node rotational degrees of freedom and a
truss or solid node none, that is not a modelling convention. It is forced by
which derivative appears in the energy. Plates and shells are the 2D version of
the same statement.
\end{mechanics}

\begin{pitfall}[why the trial space suddenly needs only \emph{one} derivative]
Compare \eqref{eq:mult} with \eqref{eq:ibp1d}. Before integrating by parts, $u$
is hit with two derivatives. After, the derivatives are shared: one on $u$, one
on $v$. Nothing was approximated --- the two lines are equal --- but the
smoothness demanded of $u$ has dropped from $C^{2}$ to ``$u'$ square
integrable.''

That trade is the entire reason the weak form exists, and it is why FEM can use
$C^{0}$ elements to solve a second-order PDE. In ME language: the principle of
virtual work never asked for the second derivative of displacement in the first
place. It only ever needed strains.
\end{pitfall}

\medskip
\noindent\textbf{Note.} Functions in $V$ are continuous.
\begin{enumerate}[label=(\arabic*),leftmargin=2em,itemsep=3pt]
  \item $\displaystyle\int_{0}^{1} v^{2}\dx < \infty$ \emph{admits}
        discontinuous functions:
        \begin{center}
        \begin{tikzpicture}[scale=0.9]
          \draw[-] (-0.3,0) -- (4.3,0);
          \draw[thick] (1.0,0) -- (1.0,0.9) -- (1.8,0.9) -- (1.8,0);
          \draw (0,-0.12) -- (0,0.12);
          \draw (4,-0.12) -- (4,0.12);
          \node[below] at (0,-0.12) {\scriptsize$0$};
          \node[below] at (4,-0.12) {\scriptsize$1$};
        \end{tikzpicture}
        \end{center}
  \item $\displaystyle\int_{0}^{1}\bigl[v'\bigr]^{2}\dx$ does \emph{not} allow
        discontinuities in $V$.
\end{enumerate}

\begin{intuition}[a jump in displacement means the bar is broken]
It is condition (2), not (1), that forces continuity --- and the mechanical
reading is immediate. If $v$ jumps at a point, the bar has separated there:
a gap has opened. The strain at that point is infinite (a finite displacement
change over zero length), so the strain energy is infinite.

$V$ therefore excludes cracked configurations automatically, without anyone
writing down a continuity requirement. Requiring finite energy \emph{implies}
$C^{0}$. Lecture 2 verifies this by explicit computation with a step function,
and the answer comes out $+\infty$ exactly as this picture predicts.

This also explains the $C^{0}$ compatibility requirement you know from element
selection: adjacent elements must share displacement along their common edge,
or the assembled field has a jump there and the energy is meaningless.
\end{intuition}

% ==================================================================
\subsection*{Summary: what the strong and weak forms each demand}

\begin{pitfall}[the weak form needs $H^{1}$, \emph{not} $C^{1}$]
An easy slip is to say ``strong needs $C^{2}$, weak needs $C^{1}$.'' The second
half is wrong, and the difference is the whole point:
\begin{align*}
  C^{1}&:\ u' \text{ exists \emph{and is continuous}}
       &&\Longrightarrow\ \textbf{no kink in } u\\[2pt]
  H^{1}&:\ \textstyle\int (u')^{2}\dx<\infty,\ u' \text{ may jump}
       &&\Longrightarrow\ \textbf{kink in } u \textbf{ allowed}
\end{align*}
$H^{1}$ is \emph{weaker} than $C^{1}$:
\[
  C^{2}\ \subset\ C^{1}\ \subset\ H^{1}\ \subset\ C^{0}.
\]
Had the weak form demanded $C^{1}$, hat functions would be illegal and
piecewise-linear FEM could not exist.
\end{pitfall}

\begin{center}\small
\renewcommand{\arraystretch}{1.25}
\begin{tabular}{@{}l l l@{}}
\toprule
 & \textbf{Strong form} ($u\in\Ck[(0,1)]{2}$) & \textbf{Weak form} ($u\in V=H^{1}_{0}$)\\
\midrule
demands on $u$   & $u,u',u''$ all continuous & $u$ continuous, $\int(u')^{2}<\infty$\\
$u'$ may jump?   & no  \xmark               & \textbf{yes} \cmark\\
$u''$ required?  & yes                       & never appears\\
equation holds   & at every point            & in an integral (averaged) sense\\
\bottomrule
\end{tabular}
\end{center}

\paragraph{What the weak form buys.}
\begin{center}\small
\renewcommand{\arraystretch}{1.2}
\begin{tabular}{@{}l l l@{}}
\toprule
 & \textbf{strong} & \textbf{weak}\\
\midrule
material interface ($a$ jumps) & no solution \xmark & fine \cmark\\
point load, re-entrant corner  & no solution \xmark & fine \cmark\\
hat functions                  & illegal \xmark     & legal \cmark \ (FEM exists only here)\\
element must compute           & $u''$              & only $u'$ (the $\bm{B}$ matrix)\\
natural BC (traction)          & imposed by hand    & free --- appears in the load vector\\
existence/uniqueness           & no theory          & Lax--Milgram ($H^{1}$ is Hilbert)\\
sparsity                       & ---                & small-support kinked basis $\Rightarrow$ banded $\Kmat$\\
\bottomrule
\end{tabular}
\end{center}

\paragraph{What it costs.}
\begin{itemize}[leftmargin=1.5em,itemsep=1pt]
  \item The PDE cannot be checked at a point --- only weighted averages of it.
  \item Strain and stress come out \emph{discontinuous} across element
        boundaries; contour plots need smoothing before presentation.
  \item Derivatives are one order less accurate than the displacement, and are
        best sampled at Gauss points rather than nodes.
  \item Conceptual overhead: Sobolev spaces and weak derivatives, i.e.\ the
        first two lectures of this course.
\end{itemize}

\begin{intuition}[the trade in one line]
\[
  \textbf{strong: } \text{demand a lot of } u \;\Rightarrow\;
  \text{pointwise information, but usually \emph{no solution}}
\]
\[
  \textbf{weak: } \text{demand little of } u \;\Rightarrow\;
  \text{a solution always exists, but only in energy terms}
\]
Engineering keeps the weak form because ``no solution at all'' is a worse
problem than ``stress needs smoothing before you plot it.''
\end{intuition}
